% Generated draft: requires mathematical and source review.
\documentclass[11pt]{article}
\usepackage[margin=25mm]{geometry}
\usepackage{fontspec}
\setmainfont{Latin Modern Roman}
\usepackage{amsmath,amssymb,mathtools,unicode-math}
\setmathfont{Latin Modern Math}
\usepackage{enumitem,needspace,xurl,hyperref}
\setlength{\parindent}{0pt}
\setlength{\parskip}{5pt}
\setlength{\emergencystretch}{3em}
\newcommand{\N}{\mathbb{N}}
\newcommand{\Z}{\mathbb{Z}}
\newcommand{\Q}{\mathbb{Q}}
\newcommand{\R}{\mathbb{R}}
\newcommand{\C}{\mathbb{C}}
\newcommand{\F}{\mathbb{F}}
\newcommand{\A}{\mathbb{A}}
\newcommand{\PP}{\mathbb{P}}
\newcommand{\E}{\mathbb{E}}
\newcommand{\eps}{\varepsilon}
\newcommand{\dd}{\,\mathrm{d}}
\begin{document}
\section*{UnsolvedMath ID: 30005874}
\textbf{Source problem code:} catalogue entry
\subsection{Definitions and mathematical statement}
Let \(S^n\) denote the smooth \(n\)-sphere. A smooth Riemannian metric \(g\) is Ricci-flat if \(\operatorname{Ric}_g=0\), and Einstein if \(\operatorname{Ric}_g=\lambda g\) for some constant \(\lambda\in\R\). A metric is called non-round here if it does not have constant sectional curvature.

The first question is:
\[
\text{Does there exist a smooth Riemannian metric }g\text{ on }S^4
\text{ such that }\operatorname{Ric}_g=0?
\]

For the second question, put
\[
D=\{(2,7),(3,6),(4,5)\},\qquad
G_{d_1,d_2}=\operatorname{SO}(d_1+1)\times\operatorname{SO}(d_2+1).
\]
For a specified smooth action \(\rho:G_{d_1,d_2}\to\operatorname{Diff}(S^{10})\), define
\[
\mathcal E_{\mathrm{nr}}(\rho)=
\left\{g:\begin{array}{l}
g\text{ is a smooth Riemannian metric on }S^{10},\\
\rho(h)^*g=g\text{ for every }h\in G_{d_1,d_2},\\
\operatorname{Ric}_g=\lambda g\text{ for some }\lambda\in\R,\\
g\text{ does not have constant sectional curvature}
\end{array}\right\}.
\]
The supplied clean statement asks, for every \((d_1,d_2)\in D\), whether the three constructed metrics exhaust this class. The supplied background describes the constructed metrics as Einstein metrics with positive scalar curvature; the clean uniqueness question imposes no positive-scalar-curvature restriction on competing metrics.

A precise uniqueness assertion requires the intended action \(\rho_{d_1,d_2}\), the three reference metrics \(g_{d_1,d_2,j}\), and an equivalence relation \(\sim\) on metrics. With those data fixed, the proposed assertion has the form
\[
\forall(d_1,d_2)\in D,\qquad
\mathcal E_{\mathrm{nr}}(\rho_{d_1,d_2})/\!\sim
=
\{[g_{d_1,d_2,1}]_{\sim},[g_{d_1,d_2,2}]_{\sim},
[g_{d_1,d_2,3}]_{\sim}\}.
\]
This display is a schematic formulation, not a fully specified conjecture: the supplied record does not identify those metrics, explicitly specify the group action, or state whether uniqueness is taken up to scaling, arbitrary diffeomorphisms, or equivariant diffeomorphisms. In particular, positive rescaling preserves the Einstein and invariance conditions, so literal equality with only three metric tensors would require an additional normalization or another qualification. The standard block action on \(S^{10}\subset\R^{d_1+1}\oplus\R^{d_2+1}\) is a possible intended action, but is not explicitly specified in the supplied text.

The original excerpt calls these the only invariant ``metrics,'' whereas the clean statement restricts the target to non-round invariant Einstein metrics. The latter restriction is used above as supplied curation. The excerpt attributes the uniqueness conjecture to numerical investigations. The two questions are recorded together; no implication between them is asserted.
\subsection{Short English statement}
The record asks whether \(S^4\) carries a smooth Ricci-flat Riemannian metric, and whether the three cited constructions exhaust the non-round invariant Einstein metrics on \(S^{10}\) for each listed pair. The second question requires clarification of the action, reference metrics, normalization, and equivalence relation. Its numerical motivation and supplied open status are retained without independent verification.
\subsection{Sources}
\begin{itemize}
\item Supplied local UnsolvedMath registry record, permanent ID \(30005874\), title: Ricci-Flat Metrics on Spheres and Invariant Einstein Metrics. Exact source problem code: \texttt{catalogue entry}. Collection: Oberwolfach Reports Open Problems.
\item \emph{Geometrie} (2024), Oberwolfach Reports, DOI 10.4171/owr/2024/28, pp. 1563--1614. Supplied source report and literature-source label: original OWR report. \url{https://doi.org/10.4171/owr/2024/28}.
\item Attribution in the supplied excerpt: Matthias Wink, \emph{Einstein metrics on Spheres}, joint work with Jan Nienhaus. The excerpt states that the uniqueness conjecture is based on numerical investigations.
\item Supplied literature-source label: Nienhaus--Wink Einstein metrics on the ten-sphere. \url{https://arxiv.org/abs/2303.04832}.
\item Literature assessment date: \texttt{2026-08-22}. This is supplied dated curation, not an independently verified assessment. It labels the Ricci-flat existence question open, describes the three non-round positive-scalar-curvature Einstein constructions on \(S^{10}\) as proved, and describes their asserted uniqueness in the indicated invariant class as conjectural. The supplied background requests independent verification of time-sensitive status and cited claims.
\item The registry labels the statement \texttt{corrected\_verified} and reports comparison with the official Oberwolfach report. These are supplied source-verification claims; no source visit or independent verification is asserted in this draft.
\end{itemize}
\end{document}
